\newpage
\subsection{\SpaceInvadersName}

Two clusterings are examined for \SpaceInvadersName, obtained at layer 4 with $K=4$ and at layer 2 with $K=5$.
At the deeper layer, three clusters track the coarse progression of the episode, from the initial, intact formation of invaders to the sparse, late-game configurations, while the fourth isolates a more specific behavior in which the agent has cleared the leftmost column and is positioned to catch the flying saucer that occasionally crosses the top of the screen.
At the shallower layer, the clustering still tracks the progression of the episode through the density of the invader formation, and separately re-isolates the cleared-left-column configuration associated with the flying saucer as its own cluster, indicating that this specific behavior remains a salient feature of the \gls{ls} across layers.
The finer clustering at layer 2 achieves a lower \gls{acu} than the coarser one at layer 4.

\begin{table}[H]
    \centering
    \begin{tblr}{
        colspec={|c|c|c|c|c|},
        hlines,
        vlines,
        row{1} = {font=\bfseries, halign=c, valign=m}
    }
        Layer & $K$ & \gls{cu}       & \gls{as}       & \gls{acu}      \\
        4     & 4   & 0.92 $\pm$0.01 & 0.69 $\pm$0.00 & 0.61 $\pm$0.00 \\
        2     & 5   & 0.92 $\pm$0.02 & 0.49 $\pm$0.01 & 0.41 $\pm$0.01 \\
    \end{tblr}
    \caption{\gls{cu}, \gls{as}, and \gls{acu} for the two clusterings analyzed in the \SpaceInvadersName environment.}
    \label{tab:qualitative_space_invaders_metrics}
\end{table}

\begin{figure}[H]
    \qualitativeClusterFigureSetup
    \qualitativeClusterTwo
        {figures/clustering_agent_representation/pictures/space_invaders/qualitative/cluster_3/image_1.png}
        {figures/clustering_agent_representation/pictures/space_invaders/qualitative/cluster_3/image_2.png}
        {Start of the episode.}
    \qualitativeClusterTwo
        {figures/clustering_agent_representation/pictures/space_invaders/qualitative/cluster_1/image_1.png}
        {figures/clustering_agent_representation/pictures/space_invaders/qualitative/cluster_1/image_2.png}
        {Middle of the episode.}
    \qualitativeClusterTwo
        {figures/clustering_agent_representation/pictures/space_invaders/qualitative/cluster_0/image_1.png}
        {figures/clustering_agent_representation/pictures/space_invaders/qualitative/cluster_0/image_2.png}
        {The left column has been cleared; the flying saucer occasionally appears.}
    \qualitativeClusterTwo
        {figures/clustering_agent_representation/pictures/space_invaders/qualitative/cluster_2/image_1.png}
        {figures/clustering_agent_representation/pictures/space_invaders/qualitative/cluster_2/image_2.png}
        {Late in the episode, few invaders remain.}
    \caption{Clusters identified by \gls{car} in the \SpaceInvadersName environment (layer 4, $K=4$).}
    \label{fig:qualitative_space_invaders}
\end{figure}


\begin{figure}[H]
    \qualitativeClusterFigureSetup
    \qualitativeClusterOne
        {figures/clustering_agent_representation/pictures/space_invaders/qualitative_alt/cluster_0/image_1.png}
        {Cluster 1: formation nearly intact.}
    \qualitativeClusterOne
        {figures/clustering_agent_representation/pictures/space_invaders/qualitative_alt/cluster_1/image_1.png}
        {Cluster 2: formation partially thinned.}
    \qualitativeClusterOne
        {figures/clustering_agent_representation/pictures/space_invaders/qualitative_alt/cluster_2/image_1.png}
        {Cluster 3: formation reorganized, three lives left.}
    \qualitativeClusterOne
        {figures/clustering_agent_representation/pictures/space_invaders/qualitative_alt/cluster_3/image_1.png}
        {Cluster 4: left column cleared, flying saucer visible.}
    \qualitativeClusterOne
        {figures/clustering_agent_representation/pictures/space_invaders/qualitative_alt/cluster_4/image_1.png}
        {Cluster 5: formation reduced to two columns.}
    \caption{Clusters identified by \gls{car} in the \SpaceInvadersName environment (layer 2, $K=5$).}
    \label{fig:qualitative_space_invaders_alt}
\end{figure}

